
\section{Architecture Overview}
The system is designed as a modular, decoupled data engineering pipeline that separates high-performance geospatial extraction from lightweight, interactive user rendering. This architecture is divided into three distinct layers:
\begin{enumerate}
  \item \textbf{Ingestion \& Processing Layer (Python/GEE/CDSE)}: Fetches and clean-masks raw multi-spectral and radar rasters. Computes spectral indices and fits unsupervised machine learning anomaly models.
  \item \textbf{Storage \& Serialization Layer (JSON/GeoJSON)}: Serializes processed anomaly data, spatial polygon boundaries, and temporal statistics into standard flat file formats.
  \item \textbf{Visualization Dashboard Layer (HTML5/CSS3/Vanilla JS)}: A futuristic single-page application (SPA) that acts as the user interface (UI) and Decision Support System (DSS), loading static serialized feeds from GitHub Pages.

\end{enumerate}
\texttt{}`
+-----------------------------------------------------------------------------------+
\begin{table}[h]
\centering
\begin{tabular}{l}
\toprule
INGESTION \& GEOSPATIAL PROCESSING LAYER \\
\midrule
+---------------------------+              +----------------------------------+ \\
+-------------+-------------+              +-----------------+----------------+ \\
\bottomrule
\end{tabular}
\end{table}+----------------|----------------------------------------------|-------------------+
                 |                                              |
                 | (Downloads Spatial Pixels \& Time Series)     |
                 v                                              v
+-----------------------------------------------------------------------------------+
\begin{table}[h]
\centering
\begin{tabular}{l}
\toprule
PYTHON PIPELINE EXECUTION ENGINE \\
\midrule
- Ingestion (ingestion.py)           - Preprocessing \& Cloud-Masking (prep.py) \\
- Spectral Indices (utils.py)        - Outlier Classification (models.py) \\
\bottomrule
\end{tabular}
\end{table}+-----------------------------------------------+-----------------------------------+
                                                |
                                                | (Serializes flat files)
                                                v
+-----------------------------------------------------------------------------------+
\begin{table}[h]
\centering
\begin{tabular}{l}
\toprule
STORAGE \& SERIALIZATION WAREHOUSE \\
\midrule
+----------------------+            +-----------------------+ \\
+----------+-----------+            +-----------+-----------+ \\
\bottomrule
\end{tabular}
\end{table}+---------------------|------------------------------------|------------------------+
                      |                                    |
                      | (Deploys statically via CI/CD)     |
                      v                                    v
+-----------------------------------------------------------------------------------+
\begin{table}[h]
\centering
\begin{tabular}{l}
\toprule
DECISION SUPPORT SYSTEM (FRONTEND SPA) \\
\midrule
- Leaflet Map (Map Canvas \& Pulsing SVG Markers) \\
- Comparison Slider (Before/After Image Clipping Canvas) \\
- Chart.js Panel (Spectral Signature Radar Charts) \\
\bottomrule
\end{tabular}
\end{table}+-----------------------------------------------------------------------------------+
\texttt{}`

By decoupling the high-performance back-end from the front-end visualization, the system bypasses the need to run resource-heavy geospatial libraries (such as GDAL, Rasterio, or Fiona) within the user's browser. Instead, the heavy computations are performed asynchronously, and the client browser only loads lightweight, pre-processed vector datasets.

---

\section{Data Ingestion Layer \& API Contracts}
The ingestion layer implements programmatic connections with three cloud registries to harvest spatial and economic datasets:

\subsection{Google Earth Engine REST API Contract}
The pipeline initializes connection with Google Earth Engine using OAuth 2.0 credentials. The Python handshake is executed as follows:
\texttt{}`python
import ee

try:
    # Initialize the GEE client under a specific GCP project billing quota
    ee.Initialize(project='tesis-500804')
    print("Google Earth Engine initialized successfully.")
except Exception as e:
    print(f"GEE Initialization failed: {e}. Initiating local mock fallback.")
\texttt{}`
When querying satellite image collections (e.g., Landsat-8/9 OLI), spatial queries are constrained using bounding boxes represented as \texttt{ee.Geometry.Polygon}. The temporal filter is applied using \texttt{ee.Filter.date(start\_date, end\_date)}.

\subsection{Copernicus CDSE REST API Contract}
To query high-resolution Sentinel-2 MSI data, the pipeline communicates with the CDSE OData endpoint.
\begin{enumerate}
  \item \textbf{Authentication}: The pipeline sends a POST request with the client credentials payload to:
    \texttt{https://identity.dataspace.copernicus.eu/auth/realms/CDSE/protocol/openid-connect/token}
  \item \textbf{OData Search Request}: A GET request queries available granules using filter strings:
    \texttt{https://catalogue.dataspace.copernicus.eu/odata/v1/Products?$filter=Attributes/OData.CSC.DoubleAttribute/any(att:att/Name eq 'cloudCover' and att/Value lt 20.0) and OData.CSC.Intersects(area=geography'SRID=4326;POLYGON((...))')}
  \item \textbf{Process API}: Fetches orthorectified surface reflectance bands directly as NumPy arrays, limiting network overhead.

\end{enumerate}
\subsection{World Bank Data REST API Contract}
Economic indicator datasets are downloaded using direct HTTP requests:
\begin{itemize}
  \item \textbf{Target URL}: \texttt{http://api.worldbank.org/v2/country/CZE/indicator/NY.GDP.MKTP.KD.ZG?format=json}
  \item \textbf{Response Processing}: The JSON payload is parsed using the Pandas library to isolate the target annual GDP growth figures, aligning the dataset with regional night-time lights anomalies.

\end{itemize}
---

\section{Processing, Masking, \& Algorithmic Pipeline}
Once raw rasters are ingested into memory as NumPy arrays, they pass through a three-stage preprocessing and model execution pipeline:

\subsection{Cloud-Masking and QA Bit-Shifting}
To eliminate cloud obstruction from spectral calculations, the pipeline performs bitwise masking operations:
\begin{itemize}
  \item \textbf{Sentinel-2 MSI}: The \texttt{QA60} band contains cloud and cirrus indicators in bits 10 and 11. The mask is calculated by shifting bits and evaluating whether both cloud flags are zero:
    \texttt{}`python
    def mask\_clouds\_s2(qa\_band):
        # Bit 10: Opaque clouds, Bit 11: Cirrus clouds
        cloud\_bit\_mask = 1 << 10
        cirrus\_bit\_mask = 1 << 11
        
        # Binary flag must be 0 for both bits to indicate clear conditions
        mask = (qa\_band \& cloud\_bit\_mask == 0) \& (qa\_band \& cirrus\_bit\_mask == 0)
        return mask
    \texttt{}`
  \item \textbf{Landsat-8/9 OLI}: The \texttt{QA\_PIXEL} band contains cloud (bit 3) and cloud shadow (bit 4) flags. The pipeline shifts these bits to isolate clear surface pixels.

\end{itemize}
\subsection{Unsupervised Anomaly Engine (Isolation Forest)}
For multi-spectral anomaly detection (e.g., Peru gold mining excavations and Chile lithium evaporation pond expansions), the processed spectral feature bands are flattened into a 2D matrix:
\begin{enumerate}
  \item \textbf{Feature Assembly}: A matrix $X \in \mathbb{R}^{N \times D}$ is constructed, where each row represents a pixel, and the columns contain computed index values ($\text{NDVI}, \text{NDWI}, \text{BSI}$).
  \item \textbf{Model Training}: An \texttt{IsolationForest} model (from \texttt{scikit-learn}) is trained with a locked random seed and configured contamination threshold:
    \texttt{}`python
    from sklearn.ensemble import IsolationForest

    # Initialize Isolation Forest with fixed seed for reproducibility
    clf = IsolationForest(
        n\_estimators=100,
        contamination=0.08,  # Expected 8\% outlier rate
        random\_state=42
    )
    
    # Fit model and predict outlier flags (-1 indicates anomaly)
    predictions = clf.fit\_predict(X)
    \texttt{}`
  \item \textbf{Spatial Projection}: Outliers are reshaped back to the original raster grid. Spatial clusters (groups of contiguous outlier pixels) are vectorized into centroid coordinates and assigned a confidence score based on their isolation path lengths.

\end{enumerate}
---

\section{Storage \& Data Warehouse Serialization}
To provide near-instantaneous page load speeds, the pipeline avoids dynamic database queries by caching processed anomalies into static files:

\subsection{anomalies.json Schema}
The serialized output contains structured metadata for each flagged anomaly event:
\texttt{}`json
[
  {
    "id": "anomaly\_madre\_de\_dios\_01",
    "type": "Deforestation/Gold Mining",
    "region": "Madre de Dios, Peru",
    "coordinates": [-12.894, -69.912],
    "confidence": 0.91,
    "date": "2026-05-14",
    "spectral\_profile": {
      "ndvi": -0.52,
      "ndwi": -0.15,
      "bsi": 0.38
    },
    "verification": {
      "ground\_truth": "Artisanal mining deforestation corridor (ASGM expansion)",
      "status": "Audited"
    },
    "details": "Unsupervised Isolation Forest flagged a sharp drop in NDVI (-0.52) and a corresponding increase in BSI (+0.38). Spatial clusters align with illegal mining corridors near the Interoceanic Highway."
  }
]
\texttt{}`

\subsection{regions.geojson Schema}
Study area boundary polygons are serialized according to the RFC 7946 GeoJSON standard:
\texttt{}`json
{
  "type": "FeatureCollection",
  "features": [
    {
      "type": "Feature",
      "properties": {
        "name": "Madre de Dios",
        "description": "Gold mining deforestation zone, Peru"
      },
      "geometry": {
        "type": "Polygon",
        "coordinates": [
          [
            [-70.5, -13.2],
            [-69.5, -13.2],
            [-69.5, -12.5],
            [-70.5, -12.5],
            [-70.5, -13.2]
          ]
        ]
      }
    }
  ]
}
\texttt{}`

---

\section{Frontend Decision Support System (DSS)}
The user interface is built as a responsive single-page dashboard using custom CSS styles and layout rules:

\subsection{Interactive Leaflet Map}
The geographic visualizer is built using \texttt{Leaflet.js} mapped to a dark tile layer:
\begin{itemize}
  \item \textbf{Basemap URL}: \texttt{https://{s}.basemaps.cartocdn.com/dark\_all/{z}/{x}/{y}{r}.png}
  \item \textbf{Pulsing SVG Markers}: Created using custom HTML \texttt{divIcon} elements. The pulsing animation uses CSS keyframes to expand concentric rings around each marker:
    \texttt{}`css
    .pulsing-marker {
        width: 12px;
        height: 12px;
        background: #00d4ff;
        border-radius: 50\%;
        position: relative;
    }
    .pulsing-marker::after {
        content: '';
        position: absolute;
        top: -6px; left: -6px;
        width: 24px; height: 24px;
        border: 2px solid #00d4ff;
        border-radius: 50\%;
        animation: pulse 1.5s infinite ease-out;
    }
    @keyframes pulse {
        0\% { transform: scale(0.5); opacity: 1; }
        100\% { transform: scale(2.0); opacity: 0; }
    }
    \texttt{}`

\end{itemize}
\subsection{Dual-Canvas Image Split Slider}
To allow users to visually inspect spatial changes without downloading large GIS files, the dashboard uses a dual-canvas layout:
\begin{enumerate}
  \item \textbf{Structure}: Two HTML5 \texttt{<canvas>} elements are layered on top of each other inside a relative-positioned container. The baseline image is drawn on the bottom canvas, and the anomaly image is drawn on the top canvas.
  \item \textbf{Clipping Event Loop}: A vertical separator bar tracks dragging inputs (\texttt{mousedown} / \texttt{touchstart}). The horizontal split percentage is calculated, and the top canvas is cropped using the CSS \texttt{clip-path} property:
    \texttt{}`javascript
    function applyClipEffect() {
        const containerWidth = container.clientWidth;
        const splitX = containerWidth * sliderSplitPercent;
        
        // Position handle element
        handle.style.left = \texttt{${sliderSplitPercent * 100}\%};
        
        // Clip top canvas (anomaly image) on the left side to reveal the baseline image
        anomalyCanvas.style.clipPath = \texttt{inset(0 0 0 ${splitX}px)};
    }
    \texttt{}`
  \item \textbf{Performance Optimization}: By using hardware-accelerated CSS clipping instead of redrawing pixel blocks on every frame, the interface achieves a smooth rendering rate of 60 frames per second (FPS), even on mobile web browsers.

\end{enumerate}
\subsection{Chart.js Spectral Radar Chart}
A radar chart displays the multi-spectral signature of the selected anomaly:
\begin{itemize}
  \item \textbf{Axes}: NDVI, NDWI, BSI, and SWIR reflectance.
  \item \textbf{Visual Styling}: The selected anomaly's signature is drawn in fluorescent cyan, and the baseline normal profile is rendered in deep blue, allowing users to quickly assess land cover anomalies.

\end{itemize}
---

\section{Deployment Architecture \& CI/CD}
The project uses GitHub Actions for deployment and scheduling:
\begin{itemize}
  \item \textbf{Build Pipeline (\texttt{test.yml})}: Installs Python dependencies, executes tests, and verifies that the code builds correctly on every pull request to \texttt{main}.
  \item \textbf{Deployment Pipeline (\texttt{deploy.yml})}: Deploys the static assets in the \texttt{/docs} directory to GitHub Pages upon merging changes.
  \item \textbf{Scheduled Pipeline (\texttt{data-refresh.yml})}: Runs weekly via a cron schedule. It executes the Python data ingestion pipeline headlessly, queries Earth Engine for new data, updates \texttt{anomalies.json}, and commits the updated cache files back to the repository.

\end{itemize}
---

\section{Security \& Vulnerability Auditing}
To secure the serverless system, the following protocols are implemented:
\begin{itemize}
  \item \textbf{API Credential Isolation}: No credentials or private keys are stored in the repository. Local testing relies on \texttt{.env} files (excluded from version control via \texttt{.gitignore}), while automated scripts access secrets configured in the GitHub Actions environment.
  \item \textbf{Content Security Policy (CSP)}: To protect against cross-site scripting (XSS) and code injection, a meta tag restricts scripts and stylesheet sources to the local domain and verified CDN hosts.
  \item \textbf{Dependency Auditing}: The CI/CD pipeline runs \texttt{pip-audit} to scan Python packages for vulnerabilities and executes \texttt{npm audit} to check frontend dependencies, ensuring the code remains secure and up to date.

\end{itemize}